\section{Classification des cycles arithmétiques irréductibles du réseau radical nonadique}

Dans la réalisation arithmétique, un nombre premier correspond à un cycle
irréductible qui ne se factorise pas en cycles positifs plus petits. La périodicité
nonadique peut sélectionner des candidats et organiser des classes résiduelles ; la
primalité reste définie par la propriété multiplicative exacte

\[
p=ab\Longrightarrow a=1\ \text{ou}\ b=1.
\]

Une fois le coproduit des cycles arithmétiques premiers sélectionné, on
obtient \cref{eq:meissel-mertens,eq:twin-prime-constant}. La structure
hélicoïdale et la racine numérique paramètrent la phase et la retenue, mais ne remplacent pas
le critère de factorisation.

\section{Densité, terme régularisé et holonomie des lacunes}

La densité positive de base est

\begin{equation}
\varepsilon_{\rm vac}=\log_{10}\frac{10}{9}.
\label{eq:vacancy-density}
\end{equation}

La séquence se déduit de la complétion
\(1000=729+271\) détermine
\begin{equation}
\varepsilon_{\rm vac}
=1-\log_{1000}(729)
=\log_{10}\frac{10}{9},
\label{eq:vacancy-density-two-forms}
\end{equation}
et le mot mécanique exact
\begin{equation}
z_n=\lfloor(n+1)\varepsilon_{\rm vac}\rfloor
-\lfloor n\varepsilon_{\rm vac}\rfloor,
\qquad n\ge1.
\label{eq:vacancy-mechanical-word}
\end{equation}
Comme \(\varepsilon_{\rm vac}\) est irrationnel, le mot est sturmien :
il a une complexité \(p(N)=N+1\), un équilibre un et une entropie topologique nulle.
Cette dynamique diffère du décalage complet sur l'alphabet
triadique : la lacune détermine une frontière ordonnée d'entropie nulle.

La somme télescopique produit
\begin{equation}
\sum_{n=1}^{N}z_n
=\lfloor(N+1)\varepsilon_{\rm vac}\rfloor,
\qquad
\sum_{n=1}^{N}(z_n-\varepsilon_{\rm vac})
=\varepsilon_{\rm vac}
-\{(N+1)\varepsilon_{\rm vac}\}.
\label{eq:vacancy-discrepancy}
\end{equation}
Les sommes partielles de la discrépance ont un module inférieur à un. C'est
l'estimation qui soutient les constantes harmoniques ultérieures.

La fraction continue commence
\[
\varepsilon_{\rm vac}
=[0;21,1,5,1,6,2,3,\ldots].
\]
C'est pourquoi les uns de \(z_n\) sont séparés par \(21\) ou \(22\) pas.
Si
\[
\sigma=22-\varepsilon_{\rm vac}^{-1}
=1-T_G(\varepsilon_{\rm vac}),
\]
alors \(a_1(\sigma)=6\) et les retours de la séparation courte sont de
longueur \(6\) ou \(7\). Les couples \(21/22\) et \(6/7\) sont deux échelles d'une
même rotation, et non des coïncidences d'entiers.

Pour le mot exact \cref{eq:vacancy-mechanical-word}, la constante régularisée est

\begin{equation}
\gamma_{\rm vac}^{+}
=\lim_{N\to\infty}\left(
\sum_{n\le N}\frac{z_n}{n}
-\varepsilon_{\rm vac}\log N
\right).
\label{eq:vacancy-gamma}
\end{equation}

L'holonomie associée est

\begin{equation}
\rho_\varepsilon(k)=\exp(2\pi\ii k\varepsilon_{\rm vac}).
\label{eq:vacancy-holonomy}
\end{equation}

Les trois quantités ont des types distincts : densité, terme fini et phase. La tour de Stieltjes des lacunes prolonge \cref{eq:vacancy-gamma} :

\begin{equation}
\gamma_{\rm vac,j}^{+}
=\varepsilon_{\rm vac}\gamma_j
+\sum_{n\ge1}\frac{(z_n-\varepsilon_{\rm vac})(\log n)^j}{n}.
\label{eq:vacancy-stieltjes}
\end{equation}

Pour \(j=0\), on retrouve \(\gamma_{\rm vac}^{+}\). La convergence exige la borne de discrépance du développement intégral des lacunes.

La convergence est prouvée ici. Pour \(j\ge0\), la suite
\(b_n=(\log n)^j/n\) tend vers zéro et est décroissante à partir de
\(n>e^j\). La borne uniforme \cref{eq:vacancy-discrepancy} et le critère de
Dirichlet prouvent la convergence de
\[
\sum_{n\ge1}(z_n-\varepsilon_{\rm vac})b_n.
\]
Pour \(j=0\), séparer
\(\varepsilon_{\rm vac}\sum_{n\le N}n^{-1}\) donne exactement
\cref{eq:vacancy-gamma}. Une sommation d'Abel borne de manière conservative la
queue après \(N\) par \(2/(N+1)\) : un terme provient du bord et
un autre de la variation totale de \(1/n\). L'exécution exhaustive déclarée utilise
\(N=10^8\), énumère les positions de Beatty des uns, vérifie les
ensembles d'écarts \(\{21,22\}\) et de retours \(\{6,7\}\), et obtient comme
partie finie
\[
-0.11207107505876366224207105242526397929\ldots.
\]
En combinant la borne analytique avec les gardes d'arrondi et de frontière
enregistrées, le certificat fini encadre
\begin{equation}
-0.112071094143613634
<\gamma_{\rm vac}^{+}
<-0.112071055516338732.
\label{eq:vacancy-gamma-interval}
\end{equation}
Cet intervalle est une valeur terminale avec erreur contrôlée. Il ne définit pas le
mot ; le mot et sa discrépance le précèdent.

Une seconde évaluation par arithmétique multiprécision et contrôle dirigé
de l'arrondi fournit, à dix-huit décimales, l'intervalle certifié
\[
-0.112071086058763562
<\gamma_{\rm vac}^{+}
<-0.112071064058763762,
\]
strictement contenu dans \cref{eq:vacancy-gamma-interval}. L'emboîtement des
deux intervalles vérifie la stabilité de la partie finie par rapport à la
troncature, aux gardes de frontière et au régime de précision. Le
résultat dépend uniquement des récurrences, des bornes de queue et de
l'arithmétique d'intervalles exposées dans ce chapitre.

\section{Contenu opératoriel de la représentation hélicoïdale}

La représentation hélicoïdale code une phase périodique accompagnée d'un
incrément de mémoire. Elle possède un contenu opératoriel lorsqu'il existe un cocycle

\[
c(\gamma\eta)=c(\gamma)+c(\eta)
\]

qui transporte la retenue et dont la monodromie s'entrelace avec l'opérateur
arithmétique. En l'absence de cette application, l'hélice ne constitue qu'une
paramétrisation de phase ; avec l'entrelaceur, elle distingue le retour de phase
du retour d'état.

\begin{remark}[Contrôle négatif]
La cofinalité des cylindres n'implique pas la conjugaison dynamique. Un
contre-exemple à la commutation entre le décalage TPK et \(T_G\)
empêche l'entrelaceur, mais n'affecte pas
\cref{eq:levy,eq:gauss-entropy}. Dans le secteur des lacunes, la divergence
de \cref{eq:vacancy-stieltjes} sous la borne déclarée réfute la famille de
constantes régularisées.
\end{remark}

\subsection*{Portée logique et domaine de validité}
Cofinalité, identités de Gauss, Lévy, entropie et Lochs : \emph{Théorème interne exact dans le système transporté}. GKW et entrelaceur dynamique : \emph{Programme ouvert avec critère de réfutation}. Densité, constante et holonomie de lacune : \emph{Théorème interne exact avec contrôle de discrépance}.
